% libre-compendium
%
% Copyright (C) 2026 Richard Qin
%
% This document is dual-licensed under either:
% - GNU Free Documentation License, Version 1.3 or any later version
%   published by the Free Software Foundation; with no Invariant Sections,
%   no Front-Cover Texts, and no Back-Cover Texts (see COPYING.DOC); or
% - Creative Commons Attribution-ShareAlike 4.0 International License
%   (see LICENSE-CC-BY-SA-4.0).

% Exercises migrated from problems04.tex and solutions04.tex.
\begin{exercise}[level=basic]
设 $z=1-\sqrt3i$。求 $|z|$、一个幅角、指数形式及 $z^{-1}$。
\begin{solution}
$|z|=2$，可取幅角 $-\pi/3$，故
\[
z=2e^{-i\pi/3}.
\]
并且
\[
z^{-1}=\frac{\overline z}{|z|^2}
=\frac{1+\sqrt3i}{4}.
\]
\end{solution}
\end{exercise}

\begin{exercise}[level=basic]
求方程 $z^4=-16$ 的全部复根。
\begin{solution}
$-16=16e^{i(\pi+2k\pi)}$，故四次根为
\[
z_k=2e^{i(\pi/4+k\pi/2)},
\qquad k=0,1,2,3.
\]
\end{solution}
\end{exercise}

\begin{exercise}[level=basic]
令 $\omega=e^{-2\pi i/5}$。计算
\[
1+\omega^2+\omega^4+\omega^6+\omega^8.
\]
\begin{solution}
指数按模 $5$ 化简后，和为
\[
1+\omega^2+\omega^4+\omega+\omega^3
=1+\omega+\omega^2+\omega^3+\omega^4=0.
\]
\end{solution}
\end{exercise}

\begin{exercise}[level=basic]
设
\[
A=\begin{pmatrix}1+i&2\\-i&3-i\end{pmatrix}.
\]
求 $A^{\mathsf H}$，并验证 $(A^{\mathsf H})^{\mathsf H}=A$。
\begin{solution}
\[
A^{\mathsf H}
=\begin{pmatrix}1-i&i\\2&3+i\end{pmatrix}.
\]
再次共轭转置，元素共轭两次并转置两次，恢复 $A$。
\end{solution}
\end{exercise}

\begin{exercise}[level=basic]
在 $\C^2$ 的标准内积下，计算
\[
x=(1+i,2)^{\mathsf T},\qquad y=(i,1-i)^{\mathsf T}
\]
的 $\langle x,y\rangle$、$\lVert x\rVert_2$，并判断是否正交。
\begin{solution}
\[
\langle x,y\rangle
=(1-i)i+2(1-i)=1+i+2-2i=3-i.
\]
且
\[
\lVert x\rVert_2=\sqrt{|1+i|^2+|2|^2}=\sqrt6.
\]
内积不为零，所以不正交。
\end{solution}
\end{exercise}

\begin{exercise}[level=basic]
判断矩阵
\[
A=\begin{pmatrix}2&i\\-i&3\end{pmatrix},
\qquad
B=\begin{pmatrix}0&i\\i&0\end{pmatrix}
\]
是否为 Hermite 矩阵。
\begin{solution}
$A^{\mathsf H}=A$，所以 $A$ 为 Hermite 矩阵。$B^{\mathsf H}=-B\ne B$，所以 $B$ 不是 Hermite 矩阵，而是斜 Hermite 矩阵。
\end{solution}
\end{exercise}

\begin{exercise}[level=basic]
把复方程组
\[
\begin{pmatrix}1&i\\2-i&1\end{pmatrix}x
=\begin{pmatrix}1+i\\2\end{pmatrix}
\]
写成四元实线性方程组。
\begin{solution}
写 $A=A_1+iA_2$，其中
\[
A_1=\begin{pmatrix}1&0\\2&1\end{pmatrix},
\qquad
A_2=\begin{pmatrix}0&1\\-1&0\end{pmatrix}.
\]
若 $x=x_1+ix_2$，则实数化方程为
\[
\begin{pmatrix}
1&0&0&-1\\
2&1&1&0\\
0&1&1&0\\
-1&0&2&1
\end{pmatrix}
\begin{pmatrix}\operatorname{Re}x_1\\\operatorname{Re}x_2\\
\operatorname{Im}x_1\\\operatorname{Im}x_2\end{pmatrix}
=\begin{pmatrix}1\\2\\1\\0\end{pmatrix}.
\]
\end{solution}
\end{exercise}

\begin{exercise}[level=basic]
写出未归一化 Fourier 矩阵 $F_2$ 和 $F_4$，并验证 $F_2^{\mathsf H}F_2=2I$。
\begin{solution}
\[
F_2=\begin{pmatrix}1&1\\1&-1\end{pmatrix},
\qquad
F_4=\begin{pmatrix}
1&1&1&1\\1&-i&-1&i\\1&-1&1&-1\\1&i&-1&-i
\end{pmatrix}.
\]
直接计算 $F_2^{\mathsf H}F_2=F_2^2=2I$。
\end{solution}
\end{exercise}

\begin{exercise}[level=standard]
使用复 Gram--Schmidt 方法把
\[
v_1=(1,i,0)^{\mathsf T},\qquad
v_2=(1,0,1)^{\mathsf T}
\]
正交归一化。
\begin{solution}
先取
\[
e_1=\frac1{\sqrt2}(1,i,0)^{\mathsf T}.
\]
因 $\langle e_1,v_2\rangle=1/\sqrt2$，
\[
w_2=v_2-\frac1{\sqrt2}e_1
=\left(\frac12,-\frac i2,1\right)^{\mathsf T}.
\]
$\lVert w_2\rVert^2=3/2$，故
\[
e_2=\frac1{\sqrt6}(1,-i,2)^{\mathsf T}.
\]
\end{solution}
\end{exercise}

\begin{exercise}[level=standard]
设 Hermite 型在标准基下的矩阵为
\[
G=\begin{pmatrix}2&i\\-i&1\end{pmatrix}.
\]
计算 $g(x,y)$，其中
$x=(1,i)^{\mathsf T}$、$y=(1,1)^{\mathsf T}$，并计算 $H(x)=g(x,x)$。
\begin{solution}
\[
Gy=\begin{pmatrix}2+i\\1-i\end{pmatrix},
\qquad
x^{\mathsf H}=(1,-i).
\]
所以
\[
g(x,y)=x^{\mathsf H}Gy=1.
\]
又
\[
Gx=\begin{pmatrix}1\\0\end{pmatrix},
\qquad H(x)=x^{\mathsf H}Gx=1.
\]
\end{solution}
\end{exercise}

\begin{exercise}[level=standard]
设新基满足 $\widetilde{\mathcal E}=\mathcal EP$，其中
\[
G=\begin{pmatrix}1&0\\0&2\end{pmatrix},\qquad
P=\begin{pmatrix}1&i\\0&1\end{pmatrix}.
\]
求 Hermite 型在新基下的矩阵。
\begin{solution}
\[
\widetilde G=P^{\mathsf H}GP
=\begin{pmatrix}1&0\\-i&1\end{pmatrix}
\begin{pmatrix}1&i\\0&2\end{pmatrix}
=\begin{pmatrix}1&i\\-i&3\end{pmatrix}.
\]
\end{solution}
\end{exercise}

\begin{exercise}[level=standard]
设
\[
A=\begin{pmatrix}2&1-i\\1+i&3\end{pmatrix}.
\]
求其特征值，验证特征值为实数，并求一组标准正交特征向量。
\begin{solution}
特征多项式为
\[
(\lambda-2)(\lambda-3)-|1-i|^2
=\lambda^2-5\lambda+4,
\]
故特征值为 $4,1$。对应特征向量可分别取
\[
v_4=(1-i,2)^{\mathsf T},
\qquad
v_1=(-2,1+i)^{\mathsf T}.
\]
二者内积为零，范数均为 $\sqrt6$，除以 $\sqrt6$ 得标准正交特征向量。
\end{solution}
\end{exercise}

\begin{exercise}[level=standard]
证明复方阵 $U$ 为酉矩阵，当且仅当它的列向量构成 $\C^n$ 的一组标准正交基。
\begin{solution}
设 $U=(u_1,\ldots,u_n)$。矩阵 $U^{\mathsf H}U$ 的第 $(i,j)$ 项为
$u_i^{\mathsf H}u_j=\langle u_i,u_j\rangle$。因此
$U^{\mathsf H}U=I$ 当且仅当各列标准正交。$n$ 个标准正交向量线性无关，故构成 $\C^n$ 的基。
\end{solution}
\end{exercise}

\begin{exercise}[level=standard]
判断下列矩阵是否正规，并说明能否酉相似对角化：
\[
A=\begin{pmatrix}1&1\\0&1\end{pmatrix},
\qquad
B=\begin{pmatrix}0&-1\\1&0\end{pmatrix}.
\]
\begin{solution}
对 $A$，
\[
A^{\mathsf H}A=\begin{pmatrix}1&1\\1&2\end{pmatrix},
\qquad
AA^{\mathsf H}=\begin{pmatrix}2&1\\1&1\end{pmatrix},
\]
二者不等，所以 $A$ 不正规，不能酉相似对角化。$B$ 是实正交矩阵，满足
$B^{\mathsf H}B=BB^{\mathsf H}=I$，故正规，可以酉相似对角化；其特征值为 $i,-i$。
\end{solution}
\end{exercise}

\begin{exercise}[level=standard]
判断 Hermite 矩阵
\[
A=\begin{pmatrix}2&i\\-i&2\end{pmatrix}
\]
是否正定，并给出一种 $A=B^{\mathsf H}B$ 分解。
\begin{solution}
特征值为 $3,1$，均为正，所以 $A$ 正定。可取 Cholesky 分解
\[
B=\begin{pmatrix}\sqrt2&i/\sqrt2\\0&\sqrt{3/2}\end{pmatrix},
\]
则直接计算得 $B^{\mathsf H}B=A$。
\end{solution}
\end{exercise}

\begin{exercise}[level=standard]
对序列 $x=(1,0,-1,0)^{\mathsf T}$ 计算四点 DFT，并用逆 DFT 恢复 $x$。
\begin{solution}
用 $F_4$ 相乘得
\[
y=(0,2,0,2)^{\mathsf T}.
\]
逆变换为
\[
x=\frac14F_4^{\mathsf H}y
=(1,0,-1,0)^{\mathsf T}.
\]
\end{solution}
\end{exercise}

\begin{exercise}[level=standard]
设 $x=(1,2,0)^{\mathsf T}$、$h=(1,-1,0)^{\mathsf T}$。计算三点循环卷积，并验证卷积定理。
\begin{solution}
按模 $3$ 计算：
\[
x\mathbin{\circledast}h=(1,1,-2)^{\mathsf T}.
\]
令 $\omega=e^{-2\pi i/3}$，则
\[
F_3x=(3,1+2\omega,1+2\omega^2)^{\mathsf T},
\]
\[
F_3h=(0,1-\omega,1-\omega^2)^{\mathsf T}.
\]
逐项相乘并利用 $1+\omega+\omega^2=0$，所得正是
$F_3(1,1,-2)^{\mathsf T}$。
\end{solution}
\end{exercise}

\begin{exercise}[level=advanced]
证明复极化恒等式，并用它从
$H(x)=x^{\mathsf H}Gx$ 恢复 $g(x,y)=x^{\mathsf H}Gy$。
\begin{solution}
展开
\[
H(x+y)-H(x-y)=4\operatorname{Re}g(x,y),
\]
以及
\[
H(x-iy)-H(x+iy)=4\operatorname{Im}g(x,y).
\]
两式按“实部加 $i$ 倍虚部”合并，得到
\[
g(x,y)=\frac14\bigl(H(x+y)-H(x-y)
+iH(x-iy)-iH(x+iy)\bigr).
\]
代入 $H(z)=z^{\mathsf H}Gz$，展开后交叉项恰合并为 $x^{\mathsf H}Gy$。
\end{solution}
\end{exercise}

\begin{exercise}[level=advanced]
设 $A$ 为复方阵。证明
\[
\ker A=\ker(A^{\mathsf H}A),
\qquad
\operatorname{Im}A^{\mathsf H}=(\ker A)^\perp.
\]
\begin{solution}
若 $Ax=0$，则 $A^{\mathsf H}Ax=0$。反之，若 $A^{\mathsf H}Ax=0$，则
\[
0=x^{\mathsf H}A^{\mathsf H}Ax=\lVert Ax\rVert_2^2,
\]
故 $Ax=0$。又对任意 $y$，
\[
y\perp\operatorname{Im}A^{\mathsf H}
\iff \langle A^{\mathsf H}x,y\rangle=0\ \forall x
\iff \langle x,Ay\rangle=0\ \forall x
\iff Ay=0.
\]
所以 $(\operatorname{Im}A^{\mathsf H})^\perp=\ker A$，取正交补得结论。
\end{solution}
\end{exercise}

\begin{exercise}[level=advanced]
设内积在某组基下的 Gram 矩阵为正定 Hermite 矩阵 $G$，算子 $T$ 的矩阵为 $A$。证明伴随算子的矩阵为
\[
G^{-1}A^{\mathsf H}G.
\]
\begin{solution}
该基下内积为 $\langle x,y\rangle=x^{\mathsf H}Gy$。若伴随矩阵为 $B$，则
\[
x^{\mathsf H}GAy=(Bx)^{\mathsf H}Gy
=x^{\mathsf H}B^{\mathsf H}Gy
\]
对所有 $x,y$ 成立，故 $GA=B^{\mathsf H}G$。取共轭转置得
$A^{\mathsf H}G=GB$，所以
\[
B=G^{-1}A^{\mathsf H}G.
\]
\end{solution}
\end{exercise}

\begin{exercise}[level=advanced]
证明正规矩阵 $A$ 满足
\[
\lVert Ax\rVert_2=\lVert A^{\mathsf H}x\rVert_2
\]
对所有 $x$ 成立，并推出
$\ker(A-\lambda I)=\ker(A^{\mathsf H}-\overline\lambda I)$。
\begin{solution}
正规性给出
\[
\lVert Ax\rVert_2^2=x^{\mathsf H}A^{\mathsf H}Ax
=x^{\mathsf H}AA^{\mathsf H}x
=\lVert A^{\mathsf H}x\rVert_2^2.
\]
$A-\lambda I$ 仍正规，因此
\[
\lVert(A-\lambda I)x\rVert
=\lVert(A^{\mathsf H}-\overline\lambda I)x\rVert.
\]
一边为零当且仅当另一边为零，故两核相等。
\end{solution}
\end{exercise}

\begin{exercise}[level=advanced]
证明未归一化 Fourier 矩阵满足
\[
F_n^2=nJ,
\]
其中 $Jx$ 把序列下标变为相反数，即 $(Jx)_j=x_{-j\bmod n}$。由此求 $F_n^4$。
\begin{solution}
对任意 $x$，
\begin{align*}
(F_n^2x)_j
&=\sum_{k=0}^{n-1}\omega_n^{jk}
\sum_{\ell=0}^{n-1}\omega_n^{k\ell}x_\ell\\
&=\sum_{\ell=0}^{n-1}x_\ell
\sum_{k=0}^{n-1}\omega_n^{k(j+\ell)}.
\end{align*}
内层和仅在 $j+\ell\equiv0\pmod n$ 时为 $n$，故
$F_n^2x=nJx$。又 $J^2=I$，所以
\[
F_n^4=n^2I.
\]
\end{solution}
\end{exercise}

\begin{exercise}[level=advanced]
设 $C$ 是第一列为 $c=(c_0,\ldots,c_{n-1})^{\mathsf T}$ 的循环矩阵。证明 Fourier 向量是 $C$ 的特征向量，并写出特征值。
\begin{solution}
令
\[
v^{(j)}=(1,\omega_n^{-j},\ldots,\omega_n^{-(n-1)j})^{\mathsf T}.
\]
循环移位使每一行的求和只差相位因子。直接作模 $n$ 下标代换可得
\[
Cv^{(j)}=\lambda_jv^{(j)},
\qquad
\lambda_j=\sum_{k=0}^{n-1}c_k\omega_n^{jk}=(F_nc)_j.
\]
因此 Fourier 向量构成循环矩阵的一组特征向量基。
\end{solution}
\end{exercise}

\begin{exercise}[level=advanced]
利用一次偶奇拆分计算 $x=(1,2,3,4)^{\mathsf T}$ 的四点 DFT，并明确写出两个二点 DFT 与蝶形合并过程。
\begin{solution}
偶下标序列为 $(1,3)$，奇下标序列为 $(2,4)$。两个二点 DFT 分别为
\[
y'=(4,-2),
\qquad
y''=(6,-2).
\]
$\omega_4=-i$。蝶形合并给出
\[
y_0=4+6=10,
\quad y_2=4-6=-2,
\]
\[
y_1=-2+(-i)(-2)=-2+2i,
\quad y_3=-2-(-i)(-2)=-2-2i.
\]
故四点 DFT 为
$(10,-2+2i,-2,-2-2i)^{\mathsf T}$。
\end{solution}
\end{exercise}

\begin{exercise}[level=advanced]
设 Hermite 矩阵
\[
A=\begin{pmatrix}2&i\\-i&2\end{pmatrix}.
\]
求 Rayleigh 商
\[
R_A(x)=\frac{x^{\mathsf H}Ax}{x^{\mathsf H}x}\qquad(x\ne0)
\]
的最大值和最小值，并写出取到两个极值的全部非零向量。
\begin{solution}
特征多项式为
\[
\det(\lambda I-A)
=(\lambda-2)^2-1=(\lambda-3)(\lambda-1),
\]
所以特征值为 $3$ 和 $1$。由 Rayleigh--Ritz 原理，
\[
\max_{x\ne0}R_A(x)=3,
\qquad
\min_{x\ne0}R_A(x)=1.
\]
对 $\lambda=3$，方程 $(A-3I)x=0$ 给出 $x_1=ix_2$，故最大值在
\[
x=c(i,1)^{\mathsf T}\qquad(c\in\C,\ c\ne0)
\]
时取到。对 $\lambda=1$，有 $x_1=-ix_2$，故最小值在
\[
x=c(-i,1)^{\mathsf T}\qquad(c\in\C,\ c\ne0)
\]
时取到。
\end{solution}
\end{exercise}

\begin{exercise}[level=advanced]
对矩阵
\[
A=\begin{pmatrix}1&1\\0&2\end{pmatrix}
\]
给出一个 Schur 酉三角分解 $A=UTU^{\mathsf H}$，并证明 $A$ 不能酉相似对角化。由此说明 Schur 酉三角化与酉对角化的适用范围有何不同。
\begin{solution}
矩阵 $A$ 本身已是上三角矩阵，因而可取
\[
U=I_2,
\qquad
T=A,
\qquad
A=UTU^{\mathsf H}.
\]
这已是一个 Schur 分解。但
\[
A^{\mathsf H}A
=\begin{pmatrix}1&1\\1&5\end{pmatrix},
\qquad
AA^{\mathsf H}
=\begin{pmatrix}2&2\\2&4\end{pmatrix},
\]
二者不等，所以 $A$ 不正规。若 $A$ 能酉相似对角化，则 $A=UDU^{\mathsf H}$，而对角矩阵 $D$ 正规，从而 $A$ 也必正规，矛盾。因此，每个复方阵都能 Schur 酉三角化，但只有正规矩阵才能酉对角化。
\end{solution}
\end{exercise}

\begin{exercise}[level=advanced]
求复矩阵
\[
A=\begin{pmatrix}1&i\\1&-i\end{pmatrix}
\]
的一个奇异值分解 $A=U\Sigma V^{\mathsf H}$，并验证 $U,V$ 均为酉矩阵。
\begin{solution}
先计算
\[
A^{\mathsf H}A
=\begin{pmatrix}1&1\\-i&i\end{pmatrix}
\begin{pmatrix}1&i\\1&-i\end{pmatrix}
=2I_2.
\]
因此两个奇异值都为 $\sqrt2$。可取
\[
V=I_2,
\qquad
\Sigma=\sqrt2 I_2,
\qquad
U=\frac1{\sqrt2}A
=\frac1{\sqrt2}\begin{pmatrix}1&i\\1&-i\end{pmatrix}.
\]
此时
\[
U^{\mathsf H}U=\frac12A^{\mathsf H}A=I_2,
\qquad
V^{\mathsf H}V=I_2,
\]
故 $U,V$ 均为酉矩阵，且
$U\Sigma V^{\mathsf H}=A$。
\end{solution}
\end{exercise}

\begin{exercise}[level=advanced]
求 Hermite 矩阵
\[
A=\begin{pmatrix}0&1+i\\1-i&0\end{pmatrix}
\]
的惯性指数 $(n_+,n_-,n_0)$，并判断二次型 $x^{\mathsf H}Ax$ 是正定、负定、半正定还是不定。
\begin{solution}
特征多项式为
\[
\det(\lambda I-A)
=\lambda^2-(1+i)(1-i)
=\lambda^2-2,
\]
所以特征值为 $\sqrt2$ 和 $-\sqrt2$，且没有零特征值。因此
\[
(n_+,n_-,n_0)=(1,1,0).
\]
Hermite 二次型在与正特征值对应的特征向量上取正值，在与负特征值对应的特征向量上取负值，故 $x^{\mathsf H}Ax$ 是不定的。
\end{solution}
\end{exercise}

\begin{exercise}[level=advanced]
在 $\C^3$ 的标准内积下，令
\[
e_1=\frac1{\sqrt2}(1,i,0)^{\mathsf T},\qquad
e_2=(0,0,1)^{\mathsf T},\qquad
x=(1,1,2)^{\mathsf T}.
\]
求 $x$ 在 $W=\operatorname{span}\{e_1,e_2\}$ 上的正交投影，并通过直接计算验证 Bessel 不等式以及正交分解的 Parseval 等式。
\begin{solution}
先有 $\langle e_1,e_2\rangle=0$ 且
$\lVert e_1\rVert=\lVert e_2\rVert=1$，所以 $e_1,e_2$ 是 $W$ 的标准正交基。投影系数为
\[
\langle e_1,x\rangle=\frac{1-i}{\sqrt2},
\qquad
\langle e_2,x\rangle=2.
\]
因而
\[
P_Wx
=\frac{1-i}{\sqrt2}e_1+2e_2
=\begin{pmatrix}(1-i)/2\\(1+i)/2\\2\end{pmatrix}.
\]
Bessel 不等式的左端为
\[
|\langle e_1,x\rangle|^2+|\langle e_2,x\rangle|^2
=1+4=5\le6=\lVert x\rVert^2.
\]
余向量为
\[
x-P_Wx
=\begin{pmatrix}(1+i)/2\\(1-i)/2\\0\end{pmatrix},
\qquad
\lVert x-P_Wx\rVert^2=1.
\]
又 $\lVert P_Wx\rVert^2=5$，故
\[
\lVert x\rVert^2
=\lVert P_Wx\rVert^2+\lVert x-P_Wx\rVert^2
=5+1=6,
\]
验证了正交分解的 Parseval 等式。
\end{solution}
\end{exercise}

\begin{exercise}[level=advanced]
采用未归一化 DFT，令 $\omega_4=-i$，并设
\[
x=(1,2,0,-1)^{\mathsf T},\qquad
y=F_4x=(2,1-3i,0,1+3i)^{\mathsf T}.
\]
分别求循环移位 $(T_1x)_k=x_{k-1\bmod4}$、调制
$(M_1x)_k=\omega_4^kx_k$ 和下标缩放 $(S_3x)_k=x_{3k\bmod4}$ 的 DFT，并验证移位、调制与缩放性质。
\begin{solution}
循环移位后
\[
T_1x=(-1,1,2,0)^{\mathsf T}.
\]
移位性质给出
\[
F_4(T_1x)_j=\omega_4^j y_j,
\]
因而
\[
F_4(T_1x)
=\bigl(2,-3-i,0,-3+i\bigr)^{\mathsf T}.
\]
另一方面，
\[
M_1x=(1,-2i,0,-i)^{\mathsf T}.
\]
调制性质给出
\[
F_4(M_1x)_j=y_{j+1\bmod4},
\]
故
\[
F_4(M_1x)
=\bigl(1-3i,0,1+3i,2\bigr)^{\mathsf T}.
\]
最后，$3^{-1}\equiv3\pmod4$，且
\[
S_3x=(1,-1,0,2)^{\mathsf T}.
\]
缩放性质给出
\[
F_4(S_3x)_j=y_{3j\bmod4},
\]
从而
\[
F_4(S_3x)
=\bigl(2,1+3i,0,1-3i\bigr)^{\mathsf T}.
\]
三个结果都可直接与 $F_4$ 相乘复核。
\end{solution}
\end{exercise}

\begin{exercise}[level=advanced]
设有限序列 $x=(1,2)$、$h=(1,-1,1)$。先计算它们的线性卷积，再将两序列补零到长度 $4$，验证补零后的四点循环卷积与线性卷积相同。说明为什么长度 $3$ 不足以避免混叠。
\begin{solution}
线性卷积长度为 $2+3-1=4$，逐项计算得
\[
x*h=(1,1,-1,2).
\]
补零为
\[
\widetilde x=(1,2,0,0),
\qquad
\widetilde h=(1,-1,1,0).
\]
四点循环卷积为
\begin{align*}
c_0&=1\cdot1+2\cdot0=1,\\
c_1&=1\cdot(-1)+2\cdot1=1,\\
c_2&=1\cdot1+2\cdot(-1)=-1,\\
c_3&=1\cdot0+2\cdot1=2.
\end{align*}
因此 $\widetilde x\mathbin{\circledast}\widetilde h=(1,1,-1,2)$，与线性卷积一致。

若只补到长度 $3$，线性卷积的第 $3$ 项会按模 $3$ 回卷并叠加到第 $0$ 项，得到三点循环卷积
\[
(3,1,-1),
\]
而不是完整的线性卷积。因此要避免混叠，长度至少应为 $4$。
\end{solution}
\end{exercise}

\begin{exercise}[level=advanced]
在 $n=4$ 时取 $x=(1,0,1,0)^{\mathsf T}$。求 $F_4x$，并验证离散不确定性估计
\[
|\operatorname{supp}x|\,|\operatorname{supp}(F_4x)|\ge4.
\]
另说明单位脉冲序列与非零常数序列也都取到等号。
\begin{solution}
直接计算得
\[
F_4x=(2,0,2,0)^{\mathsf T}.
\]
因而
\[
|\operatorname{supp}x|=2,
\qquad
|\operatorname{supp}(F_4x)|=2,
\]
从而乘积为 $4$，不确定性估计中等号成立。

若 $x=ce_r$ 是非零单位脉冲序列，则 $|\operatorname{supp}x|=1$，而
\[
(F_nx)_j=c\omega_n^{jr}\ne0
\]
对每个 $j$ 都成立，故 Fourier 支集大小为 $n$。若 $x=(c,\ldots,c)^{\mathsf T}$ 且 $c\ne0$，则它的支集大小为 $n$，而
\[
F_nx=(nc,0,\ldots,0)^{\mathsf T}
\]
的支集大小为 $1$。两种情形的支集大小乘积都为 $n$。
\end{solution}
\end{exercise}

